\chapter{结论与展望}\label{chap:conclusion}

本文围绕大语言模型数学推理后训练中的 signal collapse 问题，系统讨论了固定小组采样为何会在训练前后持续造成零梯度现象，以及难度感知自适应采样为何能够在有限预算下更有效地恢复训练信号。全文以 \texttt{OpenR1-Math-220k} 抽样数据与训练过程 checkpoint 切片为基础，从概率分析、预算效率、采样组织和系统实现四个层面展开研究，最终得到如下结论。贯穿全文的核心观点是：在可验证奖励的 reasoning 后训练中，采样预算不是优化器之外的外围参数，而是决定训练信号能否被看见的内部结构。

\section{主要结论}

第一，prompt 级正确率分布具有显著的长尾与分层特征，训练语料并不存在单一“平均难度”，而是同时包含大量接近全错的极难样本和一批高通过率样本。这意味着固定预算不可能同时满足不同难度样本的最优采样需求，也解释了为何统一组大小在 reasoning 场景下天然存在效率损失。

第二，signal collapse 并不是训练早期的暂时现象，而是在训练过程中以后验不同形式持续存在。训练前期，零信号主要来自困难样本形成的全错组；随着模型能力提升，后期零信号则越来越多地由简单样本形成的全对组主导。换言之，固定小组采样所造成的训练停滞不是阶段性偶然，而是由难度异质性与组内相对归一化共同诱发的结构性问题。

第三，自适应采样相较于固定预算具有明显优势。实验表明，无论采用平衡退出规则还是正样本优先退出规则，增量采样都能够在相近甚至更低的平均成本下显著降低信号缺失率。这说明在 reasoning 强化学习中，“采多少”并不是外在超参数，而是影响训练信号是否能够被看见的核心机制。

第四，几何块自适应虽然在统计效果上与固定块自适应接近，但其批次结构更加规整，更有利于同前缀多回答生成时的前缀复用与系统调度。因而，几何块方案的价值不只体现在少量成本差异上，更体现在它为后续 KV Cache 复用和在线推理流水线提供了更自然的执行接口。

第五，在当前数据与预算设定下，正样本优先退出规则较平衡退出规则表现出更好的预算效率。这表明在模型已有一定基础能力的中后期训练阶段，将简单题上节省下来的预算转移给中等难度样本，往往比继续在简单题上搜索 hard negative 更有价值。该发现进一步支持了“预算应按 prompt 难度再分配”的基本观点。

\section{本文贡献回收}

回到全文开头提出的研究目标，本文主要完成了三方面工作。其一，在统计分析层面，本文从 prompt 级概率分布角度重新刻画了 reasoning RL 中的零梯度现象，将其抽象为 signal collapse，并建立了与组大小、正确率分布直接相关的分析框架，从而把“为什么会学不到”解释为“预算如何错配到不同难度样本上”的问题。其二，在采样机制层面，本文围绕退出规则与采样块组织方式，系统比较了固定预算、自适应固定块和几何块策略，进一步说明难度感知预算分配在统计上确实更高效，而正样本优先规则在当前预算下具有更好的成本收益比。其三，在系统实现层面，本文并未将采样仅理解为训练前置步骤，而是将其与系统批次组织、前缀复用、KV Cache 友好性以及后续可扩展的温度缩放重要性校正联系起来，从而把问题扩展为“统计效率—系统效率”的联合设计问题。

因此，本文的核心价值不在于提出一个局部采样 trick，而在于给出了一种新的理解框架：在可验证奖励的 reasoning 强化学习后训练中，训练信号质量既由 reward 定义决定，也由 prompt 级预算如何分配、何时停止和如何组织执行所共同决定。只有把这些因素放在统一视角下考虑，后训练中的许多表面不稳定现象才能被更系统地解释。这也是本文希望回收的主线：所谓训练稳定性，并不只是梯度估计器的性质，还包括训练系统是否真正看见了足够多的可比较样本。

\section{未来展望}

未来工作可以沿三个方向继续展开。第一，在统计层面，可进一步在更多模型规模、更多基准数据集以及不同 reward 形态下验证本文结论，尤其需要考察步骤级奖励、长度惩罚和多模态 reasoning 场景下 signal collapse 的表现形式是否发生变化。第二，在方法层面，可将温度缩放、重要性采样修正和置信区间停止规则纳入统一框架，构造更强的难度感知预算分配器，使 prompt 级预算分配不仅依赖经验正确率，还能显式利用不确定性信息，并进一步与在线难度预测器或 rollout replay 机制结合。第三，在系统层面，应进一步将几何块采样与真实在线推理引擎结合，直接测量 wall-clock 时间、吞吐率、cache 命中率和前缀复用收益，从而验证统计收益是否能够稳定转化为工程收益。

总而言之，面向 reasoning 训练的强化学习后训练不应仅关注损失函数或优化器形式，还应将 prompt 级预算分配、退出规则与系统批次组织视为统一设计对象。本文希望说明，采样不是训练外部的辅助过程，而是训练内部的一部分信号机制；这一认识对于后续大规模 reasoning 模型的训练设计具有持续的研究价值。
